% Part: lambda-calculus
% Chapter: lambda-definability
% Section: pairs

\documentclass[../../../include/open-logic-section]{subfiles}

\begin{document}

\olfileid{lam}{ldf}{pai}
\olsection{جوڑے تے پچھلا عدد}

\begin{defn}
$M$ تے $N$ دا جوڑا، جیہڑا $\tuple{M,N}$ لکھدے آں، ایویں
متعین ہوندا اے۔
\[
\tuple{M,N} \ident \lambd[f][fMN].
\]
\end{defn}

بدیہی خیال ایہ اے کہ ایہ اوہ فنکشن اے جیہڑا اک فنکشن قبول
کردا اے تے اوہنوں جوڑے دے دوہاں رکناں اُتے لاؤندا اے۔ ایس
خیال توں ایہ جوڑا بناؤن والا فنکشن ملدا اے، جیہڑا دو ٹرم
لیندا اے تے اوہناں نوں رکھن والا جوڑا واپس دیندا اے۔
\[
  \fn{Pair} \ident \lambd[mn][\lambd[f][fmn]]
\]
جے جوڑا دِتّا ہووے، اسیں اوہدے رکن وی واپس حاصل کرنا چاہندے آں۔
ایس لئی دو رکن کڈھن والے فنکشن چاہیدے نیں، جیہڑے دلیل وجوں جوڑا
لیندے نیں تے اوہدے وچوں پہلا یا دوجا رکن واپس دیندے نیں۔
\begin{align*}
  \fn{Fst} & \ident \lambd[p][p(\lambd[mn][m])]\\
  \fn{Snd} & \ident \lambd[p][p(\lambd[mn][n])]
\end{align*}

\begin{prob}
  دسو کہ رکن کڈھن والے فنکشن $\fn{Fst}$ تے $\fn{Snd}$ کیوں
  کم کردے نیں۔
\end{prob}

ہُن جوڑیاں دی مدد نال پچھلا عدد دین والے فنکشن نوں
!!{lambda define} کر سکدے آں۔
\[
  \fn{Pred} \ident \lambd[n][\fn{Fst}(n (\lambd[p][\tuple{\fn{Snd}\, {p}, \fn{Succ}(\fn{Snd}\, {p})}]) \tuple{\num 0, \num 0})]
\]
یاد رکھو کہ $\num n\, f x$ دی تقلیل $f^{n}(x)$ تک ہوندی
اے۔ ایس صورت وچ $f$ اک فنکشن اے جیہڑا جوڑا $p$ لیندا اے
تے نواں جوڑا واپس دیندا اے۔ نویں جوڑے وچ $p$ دا دوجا رکن
تے اوس دوجے رکن دا جانشین ہوندا اے۔ $x$ جوڑا $\tuple{0,0}$
اے۔ سو نتیجہ $\tuple{0,0}$ اے جے $n=0$، تے ہور صورت وچ
$\tuple{\num{n-1}, \num n}$ اے۔ فیر $\fn{Pred}$ نتیجے دا پہلا
رکن واپس دیندا اے۔

گھٹاؤ دی تعریف $\fn{Pred}$ نوں $a$ اُتے $b$ واری لاؤن نال
ہو سکدی اے۔
\[
\fn{Sub} \ident \lambd[ab][b \fn{Pred}\, a].
\]

\end{document}
